\id{МРНТИ 65.35.03}{}
\vspace{2em}
\begin{header}
\swa{А.М.~Омаралиева, Ж.Т.~Ботбаева, А.Ж.~Хастаева, А.А.~Бектурганова, М.Т.~Агедилова, А.М.~Рахимжанова, А.К.~Аубакирова}{ҰННАН ЖАСАЛҒАН КОНДИТЕРЛІК ӨНІМДЕРГЕ ҚОЛДАНАТЫН ЖИДЕК ДАҚЫЛДАРЫНАН АЛЫНҒАН ҰНТАҚТАРДЫҢ ФИЗИКА ХИМИЯЛЫҚ ФУНКЦИОНАЛДЫҚ ҚАСИЕТТЕРІН ЗЕРТТЕУ}

А.М.~Омаралиева,
Ж.Т.~Ботбаева\envelope,
А.Ж.~Хастаева,
А.А.~Бектурганова,
М.Т.~Агедилова,
А.М.~Рахимжанова,
А.К.~Аубакирова
\end{header}

\begin{affil}
Қ. Құлажанов атындағы Казақ технология және бизнес университеті, Астана, Қазақстан

\corrauthor{Корреспондент-автор: zhanar.b.t88@gmail.com}
\end{affil}

Мақалада конвективті кептіру әдісімен алыyған жидек дақылдарының
(таңқурай мен қарақат жидектерінен) жасалған құрғақ ұнтақтардың физика
хиимялық антиоксиданттық аминқышқылдық және микроқұрылымдық
сипаттамаларын кешенді зерттеу нәтижелері ұсынылған. Жұмыстың мақсаты
ұннан жасалатын кондитерлік өнімдер технологиясында қолдану бағыттарын
негңздеу үшін олардың тағамдық, биологиялық және технологиялық
құндылығын бағалау.

Зерттеу нәтижелері бойынша таңқурай ұнтағының ақуыз мөлшері (7,21
±0,09\%) болса, ал майдың (3,26±0,04\%) жоғары мөлшерімен сипатталды,
сонымен қатар қарақат ұнтағында көмірсулардың жоғары (60,55 ±0,78\%)
болатындай мөлшері анықталды. Майда еритін және суда еритін
антиоксиданттардың жиынтық мөлшері ең жоғарғы деңгейде қарақат ұнтағында
байқалды (0,42±0,003 және 2,10±0,047мг/г) бұл оның жоғарғы биологиялық
белсенділігін байқатты.

Өнімдердің аминқышқылдық құрамын талдау зерттелген ұнтақтардың құрамында
алмастырылмайтын және алмастырылатын аминқышқылдарының бар екенін,
олардың ішінде лизин, лейцин, валин және аланиннің басым екенін
көрсетті, бұл олардың өсімдік текті ақуыздардың көзі ретінде бағалауға
мүмкіндік береді. Микроқұрылымдық зерттеулер таңқурай ұнтағының жоғары
дисперстілігімен және кеуектілігімен сипатталатынын, соның нәтижесінде
оның дымқылдануы артып, қамырда біркелкі таралатыны анықталды. Ал,
қарақат ұнтағы тығызырақ агрегатталған құрылымымен және кеуектілігімен
(20-30\%) ерекшеленеді.

Алынған нәтижелер жидек дақылдарынан алынған ұнтақтарды ұннан жасалатын
кондитерлік өнімдерді байытуға, оларыдың тағамдық және құрылымдық
сипаттамаларымен жақсартуға, сондай-ақ өнімнің органолептикалық және
құрылымдық сипаттамаларын жақсартуға арналған функционалдық
ингредиенттер ретінде пайдаланудың болашақтағы алдыңғы қатарлы әдіс
екендігін дәлелдейді.

{\bfseries Түйін сөздер:} таңқурай, қарақат, құрғақ ұнтақ, конвективтің
кептіруі, физика-химиялық қасиеттер, антиоксиданттық қасиеттер, құрғақ
ұнтақтардың микроқұрылымы.
\vspace{2em}
\begin{header}
ИССЛЕДОВАНИЕ ФИЗИКО-ХИМИЧЕСКИХ И ФУНКЦИОНАЛЬНЫХ СВОЙСТВ ПОРОШКОВ ИЗ ЯГОДНЫХ КУЛЬТУР ДЛЯ МУЧНЫХ КОНДИТЕРСКИХ ИЗДЕЛИЙ

А.М.~Омаралиева,
Ж.Т.~Ботбаева\envelope,
А.Ж.~Хастаева,
А.А.~Бектурганова,
М.Т.~Агедилова,
А.М.~Рахимжанова,
А.К.~Аубакирова
\end{header}

\begin{affil}
Казахский университет технологии и бизнеса им. К. Кулажанова, Астана, Казахстан,

e-mail: zhanar.b.t88@gmail.com
\end{affil}

В статье представлены результаты комплексного исследования
физико-химических антиоксидантных аминокислотных и микроструктурных
характеристик сухих порошков ягодных культур (малины и смородины),
полученных методом конвективной сушки. Целью работы была оценка их
пищевой, биологической и технологической ценности для обоснования их
применения в технологии производства мучных кондитерских изделий.

По результатам исследования, малиновый порошок характеризовался высоким
содержанием белка (7,21 ±0,09\%) и жира (3,26 ±0,04\%), а порошок черной
смородины --- высоким содержанием углеводов (60,55 ±0,78\%). Общее
содержание жирорастворимых и водорастворимых антиоксидантов было самым
высоким в порошке черной смородины (0,42 ±0,003 и 2,10 ±0,047 мг/г
соответственно), что указывает на его высокую биологическую активность.

Анализ аминокислотного состава продуктов показал, что исследованные
порошки содержат незаменимые и заменяемые аминокислоты, среди которых
преобладают лизин, лейцин, валин и аланин, что позволяет оценивать их
как источник растительных белков. Микроструктурные исследования
показали, что малиновый порошок характеризуется высокой дисперсией и
пористостью, в результате чего повышается его смачиваемость и
равномерное распределение в тесте. Порошок черной смородины, напротив,
отличается более плотной агрегированной структурой и пористостью
(20-30\%).

Полученные результаты доказывают, что использование порошков из ягодных
культур в качестве функциональных ингредиентов для обогащения мучных
кондитерских изделий, улучшения их питательных и структурных
характеристик, а также улучшения органолептических и структурных свойств
продукта является перспективным методом в будущем.

{\bfseries Ключевые слова:} малина, смородина, сухой порошок, конвективная
сушка, физико-химические свойства, антиоксидантные свойства,
микроструктура сухих порошков.
\vspace{1em}
\begin{header}
STUDY OF THE PHYSICOCHEMICAL AND FUNCTIONAL PROPERTIES OF BERRY POWDERS FOR FLOUR CONFECTIONERY PRODUCTS

A.M.~Omaraliуeva,
Zh.T.~Botbaуeva\envelope,
A.Zh.~Khastaуeva,
A.A.~Bekturganova,
M.T.~Agedilova,
A.M.~Rakhimzhanova,
A.K.~Aubakirova
\end{header}

\begin{affil}
K. Kulazhanov Kazakh University of Technology and Business, Astana, Kazakhstan,

e-mail: zhanar.b.t88@gmail.com
\end{affil}

This article presents the results of a comprehensive study of the
physicochemical antioxidant, amino acid, and microstructural
characteristics of dry berry powders (raspberries and currants) obtained
by convective drying. The aim of the study was to evaluate their
nutritional, biological, and technological value to justify their use in
flour confectionery production.

According to the study, raspberry powder was characterized by a high
protein content (7.21 ±0.09\%) and fat (3.26 ±0.04\%), while
blackcurrant powder was high in carbohydrates (60.55 ±0.78\%). The total
content of fat-soluble and water-soluble antioxidants was highest in
blackcurrant powder (0.42 ±0.003 and 2.10 ±0.047 mg/g, respectively),
indicating its high biological activity. An analysis of the amino acid
composition of the products revealed that the powders contained
essential and nonessential amino acids, predominantly lysine, leucine,
valine, and alanine, allowing them to be evaluated as a source of plant
proteins. Microstructural studies revealed that raspberry powder is
characterized by high dispersion and porosity, resulting in increased
wettability and uniform distribution in dough. Blackcurrant powder, in
contrast, has a denser aggregated structure and porosity (20-30\%).

These results demonstrate that the use of berry powders as functional
ingredients for enriching flour confectionery products, improving their
nutritional and structural properties, and enhancing the organoleptic
and structural properties of the product is a promising approach for the
future.

{\bfseries Key words:} raspberry, currant, dry powder, convective drying,
physicochemical properties, antioxidant properties, microstructure of
dry powders.

\begin{multicols}{2}
{\bfseries Кіріспе.} Ұннан жасалатын кондитерлік өнімдер (печенье пряниктер
және тағы басқа) қолжетімділігі, ұзақ сақталу мерзімі және тұрақты
сұранысқа ие күнделікті тұтынатын, ең кең таралған өнімдердің қатарына
жатады. Сонымен қатар, дәстүрлі рецептуралар оңай сіңетін көмірсулардың
жоғары мөлшерімен және тағамдық талшықтардың, витаминдердің және
биологиялық белсенді заттардың салыстырмалы түрде төмен үлесімен
сипатталады, бұл олардың тағамдық және биологиялық құндылығын арттыруға
бағытталған жіне тәсілдерді әзірлеу қажеттілігін туындатады {[}1,2{]}.

Соңғы жылдары, әлемдік тағамдық технологиялар практикасында жеміс-жидек
ұнтақтарын сондай-ақ, оларды өңдеуі кезінде алынатын жанама өнімдерді
функционалдық ингредиенттер ретінде пайдалану бағыты белсенді дамып
келеді. Әдебиеттерге шолу сипатындағы зерттеулер мұндай қоспалардың
тағамдық талшықтардың, полифенолдардың органикалық қышқылдардың және
дайын өнімдердің антиоксиданттық белсенділігі мен нутриенттік тығыздығын
арттыруға қабілетті табиғи пигменттердің құнды көзі болып табылатынын
көрсетеді.

Шетелдік зерттеулерде антоциондар мен басқа да фенолдық қосылыстардың
жоғарғы мөлшерімен ерекшеленетін жидек дақылдарын пайдалануға ерекше
көңіл бөлінуде. Қарақат жидегінің антиоксиданттары жоғары биологиялық
белесенділігімен ерекшеленетін алдыңғы қатарлы сапалы көздерінің бірі
ретінде қарастырылады {[}2,6,8{]}. Эксперименттік жұмыстар қарақат
ұнтақтары мен сығымдалған қалдықтарын печенье және басқа да ұннан
жасалатын өнімдер технологиясында қолданудың тиімділігін растайды, бұл
дайын өнімдердің антиоксиданттық әлеуетінің артуымен және функционалдық
қасиеттерінің жақсаруымен қатар жүреді {[}3{]}.

Ұқсас үрдістер зерттеу барысында, таңқурай ұнтақтарын қолдану кезінде де
анықталды. Таңқурай ұнтағын енгізу өнімдердегі фенолдық қосылыстардың
мөлшерін және антиоксиданттық белсенділікті арттыруға, сондай-ақ дәм иіс
профилінің айқын қалыптасуына әсер ететіні көрсетілген, алайда бұл дайын
өнімнің сынғыштығы мен құрылымдық механикалық сипаттамаларын кешенді
бағалай қажеттілігін көрсетеді {[}4{]}. Осыған байланысты шетелдік
авторлар байытылған өнімдердің химиялық құрамымен қатар
құрылымдық-механикалық сипаттамаларын кешенді бағалау қажеттігін атап
көрсетеді {[}5{]}.

Жидек ұнтақтарының сапасы мен функциоанлдық қасиеттерін айқындайтын
маңызды факторларының бірі - оларды бөліп алу әдісі. Зерттеулер
өнімдерді кептіру параметрлерінің полифенолдарының С дәруменінің,
антициандары мен каротиноидтарының сақталуына, сондай ақ бөлшектерінің
морфологиясы мен материалдың кеуектілігінің қалыптасуына елеулі әсер
ететінін көрсетеді {[}6-8{]}. Дұрыс таңдалмаған режимдерде жүргізілген
конвективті кептіру термолабильді қосылыстардың ішінара деградациясына
және ұнтақтардың құрылымдық сипаттамаларының нашарлауына әкелуі мүмкін
екені анықталған {[}9{]}.

Сонымен қатар пісіру процесі кезінде термиялық өңдеу биологиялық
белсенді заттардың қосымша азаюына себеп болуы мүмкін. Жидек
сығындыларымен байытылған ұннан жасалатын өнімдердегі полифенолдардың
тұрақтылығын зерттеу олардың сақталу дәрежесі пісірудің температуралық
уақыттық параметрлеріне және қамыр жүйесінің құрамына тәуелді екендігі
анықталды {[}10{]}. Осыған орай ұнтақтардың микроқұрылымы, олардың
химиялық құрамы және дайын өнім құрамындағы биоактивті компоненттердің
тұрақтылығы арасындағы өзара байланысты зерттеу ерекше мағызға ие.

Шетелдік жарияланымдардың едәуір санына қарамастан, зерттеулердің басым
бөлігі шикізаттың жекелеген түрлеріне немесе көрсеткіштердің жеке
топтарына бағытталған, ал жидек ұнтақтарының физика-химиялық
қасиеттерін, антиоксиданттық белсенділігін, аминқышқылдық құрамын және
морфометриялық сипаттамаларын біріктере отырып жүргізілген кешенді
жұмыстары әлі де шектеулі болып отыр, әсіресе отындық шикізатты
пайдалану жағдайында {[}11,12{]}.

Қазақстан жағдайында жергілікті жидек дақылдарын пайдалану және оларды
терең өңдеу технологияларын әзірлеу арқылы функционалдық ингредиенттер
алу өзекті мәнге ие болуда, бұл жаппай тұтынылатын тағам өнімдерінің
ассортиментін кеңейтуге және олардың тағамдық құндылығын арттыруға ықпал
етеді.

Осыған байланысты аталған зерттеудің мақсаты конвективті кептіру
әдісімен алынған жидек дақылдарынан (таңқурай мен қарақаттан) жасалған
ұнтақтардың физика- химиялық көрсеткіштерін, антиоксиданттық
белсенділігін, аминқышқылдық құрамын және микроқұрылымдық сипаттамаларын
кешенді түрде бағалау арқылы оларды ұннан жасалатын кондитерлік өнімдер
технологиясында пайдалану бағыттарын негіздеу болып табылады.

{\bfseries Материалдар мен тәсілдер.} Зерттеу нысандары ретінде 2024 және
2025 жылдардағы өнімнен алынған жидек дақылдары (таңқурай мен қарақат)
пайдаланылады. Аталған дақылдардан әртүрлі температуралық режимде, қабат
қалыңдағанда және үрдіс қиындағанда конвективті кептіру әдісі арқылы
құрғақ ұнтақтар алынды.

Жидек дақылдарынан алынған ұнтақтардың физика-химиялық қасиеттерін
зерттеу стандартты жағдайларда жүргізілді: температура - 21
\tsp{0}С, ауаның салыстырмалы ылғалдылығы - 64\%.
Көрсеткіштерді анықтау қолданыстағы нормативтік құжаттарға сәйкес жүзеге
асырылды. Ақуыздың массалық үлесі Кьелдаль әдісімен, майдың массалық
үлесі Сокслет әдісімен, көмірсулардың массалық үлесі МЕМСТ 54037-2010
стандартына сәйкес перманганатометриялық әдіспен анықталды.

Антиоксиданттардың мөлшерін анықтау ГОСТ Р54037-2010 талаптарына сәйкес
Цвет Яуза-01-АА аспабында амперметриялық әдіспен жүргізілді. Барлық
өлшеулерде стандарт ретінде кварцетин қолданылды.

Құрғақ ұнтақтардың микроқұрылымдарын талдауы 5-10х үлкейту кезінде
оптикалық микроскопия әдісімен жүргізілді. Суреттерді өңдеу үшін
бағдарламалық қамтамасыз етуі пайдаланып, бөлшектрердің орташа өлшемі,
диаметрлер бойынша таралу диапазоны, дөңгелектік коэффициенті және
ұнтақтардың салыстырмалы кеуектілігі анықталды.

Микроскопиялық суреттерді алу үшін әрбір ұнтақтардың аз мөлшері шыны
негізге біркелкі таралады. Суреттер 5х объективін қолдана отырып,
қараңғы өріс режимінде түсірілді. Фокус тереңдігін арттыру мақсатында
брекетинг әдісі қолданылды, яғни үлгінің бір аймағында әртүрлі фокуста
бірнеше сурет түсіріліп, кейін олар Helicon Focus бағдарламасында
біртұтас бейнеге біріктірілді.

Барлық эксперименттік зерттеулер кемінде үш қайталымда жасалады. Аланған
деректер Microsoft Exel пакеті мен Statistica 10.0 бағдарламасын
пайдалана отырып варияциялық статистика әдістерімен өңделді.Нәтижелер
орташа мәндер және олардың стандарттық ауырқулары (М±SD) түрінде
ұсынылды. Орташа мәндер арасындағы айырмашылықтардың статистикалық
сенімділігі р<0.05 мәнділік деңгейінде Стьюдент критерийі арқылы
бағаланады.

{\bfseries Нәтижелер мен талқылау.} Жидек дақылдарынан алынған ұнтақтардың
физика-химиялық қасиеттерін зерттеу олардың тағамдық және биологиялық
құндылығын бағалау, сондай-ақ тағам өнеркәсібінде практикалық қолдану
бағыттарын анықтау үшін маңызды болып табылады. Талдау барысында кейбір
критерийлер негізінде ақуыздың, майлардың және көмірсулародың массалық
үлесі алынды, себебі дәл осы компоненттер өнімнің тағамдық құндылығы мен
функционалдық қасиеттерін айқындайды. Көрсеткіштерді салыстырмалы түрде
зерттеу дақылдар арасындағы айырмашылықтарды анықтауға,
олардыңартықшылықтары мен кемшіліктерін белгілеуге, сондай --ақ оларды
байытқыш немесе функционалддық қоспалыр ретінде пайдалануға болады.

1-кестеден жидек дақылдарынан алынған ұнтақтардың физика-химиялық
қасиеттерінің зерттеу нәтижелерін көруге болады.

1-кестеден жидек ұнтақтарының физика-химиялық көрсеткіштерінің құрамдары
бойынша бір-бірінен ерекшеленетіндігі байқалды. Ақуыздың массалық үлесі
таңқурай ұнтағында (7,21±0,09\%) қарақат ұнтағына (6,16±0,05\%)
қарағанда жоғарырақ, бұл бірінші үлгінің ақуыздық құндылығы жоғары
екендігі дәлелденді.
\end{multicols}

\tcap{1-кесте. Жидек дақылдарынан алынған ұнтақтардың физика-химиялық көрсеткіштері (M±SD, n=3)}
\begin{longtblr}[
  label = none,
  entry = none,
]{
  width = \linewidth,
  colspec = {Q[388]Q[269]Q[283]},
  cells = {c},
  cells = {font = \small},
  row{1} = {font=\bfseries\small},
  hlines,
  vlines,
}
Берілген көрсеткіштер & Қарақат жидегінің ұнтағы & Таңқурай жидегінің ұнтағы\\
Ақуыздың анықталған массалық үлес,\% & 6,16±0,05 & 7,21±0,09\\
Майдың массалық үлесі,\% & 0,95±0,01 & 3,26±0,04\\
Көмірсулардың массалық үлесі, \% & 60,55±0,53 & 41,16±0,53
\end{longtblr}

\begin{multicols}{2}
Майдың массалық үлесі де таңқурай жидегінің ұнтағында едәуір жоғарылап,
3,26±0,04\% құраса, қарақат жидегінің ұнтағынада бұл көрсеткіш не бәрі
0,95±0,01\% - ды құрады. Сонымен таңқурай жидегінің ұнтеғы ақуыздық және
май мөдшерінің жоғарылығымен ерекшеленсе, қарақат жидегінде көмірсул
көрсеткіші жоғары болды.

Алынған деректер әрбір зерттелінген ұнтақ түрінің өзіндік тағамдық
құндылығының жоғары екендігін дәлелдеді, сонымен қатар олар тағам
өнеркәсібінің әртүрлі бағыттарында мақсатты түрде пайдалануға мүмкіндік
береді.

Құрғақ ұнтақтардағы майда еритін және суда еритін антиоксиданттардың
жиынтық мөлшерін амперметриялық әдістер арқылы анықтауға болады. Құрғақ
жидек ұнтақтардағы майда еритін және суда еритін антиоксиданттардың
жиынтық мөлшері амперметриялық әдіс арқылы анықталды. Мұндай әдістер
берілген потенциал кезінде жұмысшы электродта фенолдық табиғатты
қосылыстардың тотығуы нәтижесінде пайда болатын электр тоғын тіркеуге
негізделген және жоғары сезімталдығы мен селективтітігімен ерекшеленеді
{[}13{]}. Талдау нәтижелері 2- кестеде көрсетілген.
\end{multicols}

\tcap{2-кесте. Жидек дақылдарынан алынған құрғақ ұнтақтардағы антиоксиданттардың жиынтық мөлшері (құрғақ затқа мг/г) (М±SD, n=3)}
\begin{longtblr}[
  label = none,
  entry = none,
]{
  width = \linewidth,
  colspec = {Q[237]Q[300]Q[300]},
  cells = {c},
  cells = {font = \small},
  row{1} = {font=\bfseries\small},
  hlines,
  vlines,
}
Зерттелетен нысанның атауы & Майда еритін антиоксиданттар көрсеткіштері, мг/г & {\textbf{Судла еритін}\\\textbf{антиоксиланттардың көрсеткіштері, мг/г}}\\
Таңқурай жидегінің ұнтағы & 0,18±0,001 & 0,73±0,0044\\
Қарақат жидегінің ұнтағы & 0,42±0,003 & 2,10±0,047
\end{longtblr}

\begin{multicols}{2}
Деректерді талдау (2 - кестеде келтірілген) антиоксиданттардың ең жоғары
мөлшері қарақат ұнтағында шоғырланғанын көрсетті: майда еритін
антиоксиданттардың көрсеткіштері: 0,42±0,003 мг/г, суда еритін
антиоксиданттардың көрсеткіштері 2,10 ±0,047 мг/г.

Қарақат ұнтағындағы суда еритін антиаосиданттардың мөлшері (2,10±0,047
мг/г) таңқурай ұнтағымен салыстырғанда (0,73±0,004 мг/г) статистикалық
көрсеткіштер алынған нәтижелердің жоғары екендігін көрсетті
(р<0,05).

Қарақаттағы антиоксиданттардың жоғары концентрациясы оның құрамындағы С
витамині, катехиндер мен флаваноидтарға байланысты, бұл қосылыстар
кептіру процесңнен кейін де биологиялық белсенділігін сақтайды.
Токоферолдар мен каротиноидтарды қамтитын майда еритін антиоксиданттық
қосылыстар өнімнің функционалдық құндылығын арттырып қана қоймай, тотығу
үрдістеріне де төзімділігін күшейтеді.

Осылайша қарақат ұнтақтарын кондитерлік өнімдердің рецептураларында
пайдалануға арналған биологиялық белсенді қосылыстардың алдыңғы қатарлы
көзі ретінде қарастыруға болады. Құрғақ жидек ұнтақтарын печенье мен
пряник құрамына енгізу дайын өнімдердің тағамдық құндылылығын арттырып
қана қоймай, айқын түс беруге және органикалық қасиеттерін жақсартуға
мүмкіндік береді.

Жалпы алғанда, қарақат пен таңқурай ұнтақтары биологиялық құндылықты
арттырып қана қоймай, ұннан жасалған кондитерлік өнімдердің сапасын
жақсартуға қабілетті функционалдық ингредиенттер ретінде жоғары әлеуетке
ие.

Аминқышқылдық құрам ақуыздың биологиялық құндылығын және оның адамдардың
тамақтануында функционалдық рөлін айқындайтын маңызды көрсеткіш болып
табылады. Аминқышқылдық профильді талдау алмастырылмайтын және
алмастырылатын аминқышқылдардың сандық және сапалық айырмашылықтарын
анықтауға мүмкіндік береді, бұл композиттік қоспалар мен функционалдық
тағам өнімдерін әзірлеу барысында ерекше маңызға ие {[}14{]}. Ақуыз
құрамы шикізаттың түрі мен сортына, сондай ақ жеке аминқышлылдардың
сақталуына әсер етуі мүмкін технологиялық факторларға (кептіру,
ұнтақтау) тәуекелді екенін ескерту қажет {[}15,16{]}.

Жидек дақылдарынан алынған ұнтақтардың аминқышқылдық құрамын зерттеу
нәтижелері кестеде келтірілген.
\end{multicols}

\tcap{3-кесте. Ұнтақ үлгілерінің аминқышқылдық құрамын салыстырмалы талдау көрсеткіштері \% (М±SD, n=3)}
\begin{longtblr}[
  label = none,
  entry = none,
]{
  cells = {c},
  cells = {font = \small},
  row{1} = {font=\bfseries\small},
  hlines,
  vlines,
}
№ & Аминқышқылдар & Таңқурай жидегі & Қарақат жидегі\\
1 & Аргинин & 0,451±0,180 & 0,382±0,153\\
2 & Лизин & 0,341±0,116 & 0,310±0,106\\
3 & Тирозин & 0,239±0,072 & 0,199±0,060\\
4 & Фенилаланин & 0,359±0,108 & 0,366±0,110\\
5 & Гистидин & 0,050±0,025 & 0,045±0,022\\
6 & Лейцин+изолейцин & 0,442±0,115 & 0,398±0,103\\
7 & Метионин & 0,064±0,022 & 0,103±0,035\\
8 & Валин & 0,423±0,169 & 0,287±0,115\\
9 & Пролин & 0,396±0,103 & 0,398±0,103\\
10 & Тренонин & 0,442±0,177 & 0,366±-0,146\\
11 & Серин & 0,479±0,124 & 0,310±0,081\\
12 & Аланин & 0,571±0,148 & 0,326±0,085\\
13 & Глицин & 0,368±0,125 & 0,294±0,100
\end{longtblr}

\begin{multicols}{2}
3 кестеде келтірілген деректердің талдау нәтижелерінің негізінде
таңқұрай және қарақат жидектерінің ұнтақтарында алмастырылатын және
алмастырылмайтын аминқышқылдардың кең спектрі бар, бұл олардың
функционалдық тағам өнімдерінің компененттері ретінде әлеуетті
биологиялық құндылығын көрсетеді.

Зерттеу нәтижелері бойынша, таңқурай ұнтағын қарақат ұнтағымен
салыстырғанда зерттелген аминқышқылдарының көпшілігі бойынша жоғары
көрсеткіштерге ие. Ең айқын айырмашылықтар аргинин (0,451±0,180
көрсеткіші 0,382±0,153\% көрсеткішіне айырмашылық байқалса, лизиннің
(0,341±0,116\% көрсеткіші 0,310±0,106\%) валиннің көрсеткіші
(0,423±0,169\% қарсы 0,287±0,115\%), серин (0,479±0,124\% көрсеткішіне
0,310±0,081\% және аланин (0,571±0,148\% көрсеткішіне қарсы
0,326±0,085\% болатын құрамы байқалады. Осы амин қышқылдарынығ
мөлшерлерінің жоғары болуы таңқурай жидегінің ұнтағының ақуыздық және
метаболитикалық құндылығының жоғары екендігін көрсетеді.

Айта кету керек, лизин -- дәнді дақылдардан алынған өнімдерінде шектеулі
аминқышқылдарының бірі, ал таңқұрай ұнтағында оның концентрациясы
солыстырмалы түрде жоғары, бұл ингредиентті ұннан жасалатын өнімдерде
қолдану кезінде мағызды болып табылады. Лейцин мен изолейциннің
(0,442±0,115\%) және валиннің (0,423±0,169\%) көп мөлшерде болуы
таңқурай ұнтағының дайын өнімдердің ақуыздық құрамын толықтырудағы
әлеуетін көрсетеді.

Қарақат ұнтағы, керісінше, метионин (0,103±0,035 \%) және фенилаланин
(0,366±0,110\%) құрамында таңқурай ұнтағына қарағанда сәл жоғары. Бұл
аминқышқылдар ақуыз синтезінде, метаболизмді реттеуде және ароматикалық
қосылыстардың түзілуінде маңызды рөл атқарады, бұл дайын өнімдердің
сапасына оң әсер етуі мүмкін.

Пролиннің мөлшері екі үлгіде де шамамен бірдейде (0,396-0,398\%) бұл
ақуыз компаненттерінің құрылымдық қасиеттерінің ұқсастығын көрсетеді.
Глицин мне сериннің болуы ақуыздың функционалдық қасиеттерінің
ұқсастығын көрсететді. Глицин мен сериннің болуы ақуыздың функционалдық
қасиеттерін, соның ішінде су ұқстау қабілетін және өнімдердің
текстурасын қалыптастыруға қатысуын жақсартады.

Жалпы, алынған нәтижелер жидек ұнтақтары физиологиялық тұрғыдан маңызды
аминқышқылдарды қамтитын және дәстүрлі ұндық шикізаттың аминқышқылдық
профилін толықтыра алатынын көрсетеді, бұл дайын өнімдердің ақуыздарының
биологиялық толықтығын арттырады. Осылайша, жидек ұнтақтарын композиттік
қоспалар мен ұннан жасалатын өнімдер құрамына енгізу ақуыздардардың
аминқышқылдық спектірін кеңейтумен қатар олардың биологиялық және
функционалдық құндылығын арттыруға мүмкіндік береді.

Құрғақ ұнтақтардың микроқұрылымын талдауы -- алынған өнімнің
морфологиялық ерекшеліктерін анықтап, кептіру параметрлерінің жасушалық
құрылымның сақталуына әсерін бағалауға мүмкіндік беретін маңызды кезең.
Белгілі болғандай, бөлшектердің формасы, кеуекті құрылымды болады,
фракция көлемі және беттің сипаты ерігіштік, су ұстамдылық қабілеті,
төгілуге қабілеттілік және биологиялық белсенді қосылыстарды ұстап тұру
сияқты қасиеттермен тікелей байланысты.

1-суретте таңқурай ұнтағының микроқұрылымдары көрсетілген. Таңқурай
ұнтағы ұсақ дисперсті құрылымымен ерекшеленеді: бөлшектердің орташа
көлемі 150-300 мкмдөңгелектілік коэффициенті 0,70-0,80. Кеуектілігі өте
жоғары (45\%-дейін), бұл ұнтақты тез сулануға және жылдам еріп кетуіне
мүмкіндік береді.
\end{multicols}

\begin{figs}[1-сурет. Таңқурай ұнтағының микроқұрылымы, 5х және 10х
ұлғайтқышы]
    \fig[0.45\textwidth]{p2/image4}
    \fig[0.45\textwidth]{p2/image5}
\end{figs}



1-ші суреттен көрінетіндей, бөлшектер бос құрылымға және дөңгелек
формаға ие, бұл ұнтақұтың қамырда біркелкі таралуын қамтамасыз етеді
және ерігіштігін жақсартады.

Қарақат ұнтағы бөлшектерінің көлемі бойынша кең ауқымға ие ()200- 700
мкм) және тығыз агрегаттар түзеді. Кеуектілік 20-300 \%, дөңгелектілік
коэффициенті 0,50-0,60. Осы деректерден байқалатын негізгі ерекшелік:
таңқурай ұнтағы ұсақ, жеңіл және жақсы ергіш құрылымымен сипатталса,
қарақат ұнтағы көлемді, тығыз және, аз кеуекті құрылымға ие, бұл оның
қамырда таралуы қасиетіне және ерігіштігіне әсер етеді.

\fig{p2/image6}[]\fig{p2/image7}[]

{\bfseries 2-сурет. Қарақат ұнтағының микроқұрылымы, ұлғайтқышы 5х}

Сурет 2 ден көрініп тұрғандай, бөлшектер тығыз және біркелкі емес
құрылымға ие, кей жерлерде беттің қатаюы байқалады, бұл
карамелизациясына және антициандардың жоғары концентрациясына
байланысты.

Зертелген ұнтақтардың микроқұрылымдарының сандық сипаттамаларын жүйелеу
мақсатында бөлшектердің орташа көлемі, дөңгелектілік коэффициенті және
салыстырмалы кеуектілік құрылымы есептеледі. Нәтижелер 4- кестеде
келтірілген.

{\bfseries 4-кесте. Жидек дақылдарының құрғақ ұнтақтарының
микроқұрылымдарының морфометриялық көрсеткіштері}

%% \begin{longtable}[]{@{}
%%   >{\raggedright\arraybackslash}p{(\linewidth - 10\tabcolsep) * \real{0.1613}}
%%   >{\centering\arraybackslash}p{(\linewidth - 10\tabcolsep) * \real{0.1442}}
%%   >{\centering\arraybackslash}p{(\linewidth - 10\tabcolsep) * \real{0.1796}}
%%   >{\centering\arraybackslash}p{(\linewidth - 10\tabcolsep) * \real{0.1618}}
%%   >{\centering\arraybackslash}p{(\linewidth - 10\tabcolsep) * \real{0.1156}}
%%   >{\centering\arraybackslash}p{(\linewidth - 10\tabcolsep) * \real{0.2374}}@{}}
%% \toprule\noalign{}
%% \begin{minipage}[b]{\linewidth}\centering
%% {\bfseries Үлгілер}
%% \end{minipage} & \begin{minipage}[b]{\linewidth}\centering
%% {\bfseries Ұнтақтардың орташа өлшемдері, мкм}
%% \end{minipage} & \begin{minipage}[b]{\linewidth}\centering
%% {\bfseries Өлшемдердің диапазондары, мкм}
%% \end{minipage} & \begin{minipage}[b]{\linewidth}\centering
%% {\bfseries Дөңгелектік коэффициенті}
%% \end{minipage} & \begin{minipage}[b]{\linewidth}\centering
%% {\bfseries Кеуектілік}
%% 
%% {\bfseries \%}
%% \end{minipage} & \begin{minipage}[b]{\linewidth}\centering
%% {\bfseries Құрылымдардың ерекшелігі}
%% \end{minipage} \\
%% \midrule\noalign{}
%% \endhead
%% \bottomrule\noalign{}
%% \endlastfoot
%% Таңқурайдың құрғақ ұнтағы & 150-300 & 400 & 0,70-0,80 & 45 & Ұсақ
%% бөлшектері агрегирленген, құрылымы кеуекті \\
%% Қарақаттың құрғақ ұнтағы & 200-700 & 700 & 0,50-0,60 & 20-30 &
%% Агрегаттары тығыз,фрагменттері нығыздалған, пигментациясы өте жоғары \\
%% \end{longtable}

Морфометриялық сипаттамаларды талдау (4- кесте) келтірілді,жидет
ұнтақтарының ұсақ өлшемдері (таңқурай мен қарақат) олардың жоғары
дисперстілігімен ерекшеленіп олардың ерігіштік қасиеттерін жақсартады,
бұл фенолдық қосылыстардың сақталуын қамтамасыз ететін жоғары
кеуектілігімен ерекшеленеді.

{\bfseries Қорытындылар.} Жүргізілген зерттеулер нәтижесінде конвективті
кептіру әдісімен алынған таңқурай мен қарақат ұнтақтарының жоғары
тағамдық, биологияялық және технологиялық құндылыққа ие екені анықталды,
бұл оларды ұннан жасалатын кондитерлік өнімдер технологиясында
қолданудың орындылығын негіздейді.

Таңқурай ұнтағының ақуыз бен май мөлшерінің жоғары екені, сондай-ақ,
жоғары дисперстілігі мен кеуектілігі оның жақсы сулануын және қамыр
жүйесінде біркелкі таралуын қамтамасыз ететіні тәжірибелер негізінде
дәлелденді.

Қара қарақат ұнтағы майда еритін суда еритін антиоксидаттардың алдыңғы
қатарлы негізгі көзі болып табылатыны анықталды, олардың мөлшері
тиісінше 0,42-±0,003 және 2,10±0,047 мг/г деңгейіне жетеді, бұл
өнімдердің биологиялық құндылығы мен тотығу тұрақтылығыының артуына
ықпал етеді.

Аминқышқылдылық құрамды талдау зерттелген ұнтақтарда алмастырылатын және
алмастырылмайтын аминқышқылдарының кешені бар екенін көрсетті. Таңқурай
ұнтағы лизин, валин, ленйцин серин жәнек аланин мөлшерінің жоғары
болуымен ерекшеленеді, бұл оның жоғары ақуыздық және функционалдық
құндалаығын дәлелдейді.

Микроқұралымдық зерттеулер ұнтақтардың морфологиялық ерекшеліктері
бастапқы шикізаттың түріне және кептіру жағдайларына тәуелді екенін
көрсетті

Таңқұрай ұнтағының жоғары кеуектілігі мен ұсақ дисперсті құрылымы оның
технлогиялық қасиеттерін жақсартса, қара қарақат ұнтағының тығызырақ
құрылымы биоллогиялық белсенді қосылыстардың ықпалына әсер етеді.

Физика-химиялық антиоксидаттық, аминқышқылдық және құрылымдық
сипаттамалардың кешенді бағасы жидек дақылдары ұнтақтарын ұннан
жасалатынын кондитерлік өнімдерді байыту және олардың тағамдық әрі
биологиялық құндылығын арттыру мақсатында функционалдық ингредиентитер
ретәінде қолданудың болашағы жатыр.

Алынған нәтижелер бойынша ұсақ ұнтақтардың тұтынушылық қасиеттері
жақсартылған және функционалдық бағыттағы печенье мен пряник өнімдерінің
жаңа рецептураларын әзірлеу кезінде пайдалануы мүмкін.

\emph{{\bfseries Қаржыландыру}. Бұл зерттеуді Қазақстан Республикасы Ғылым
және жоғары минситрлігінің Ғылым комитеті тарапынан IRN №BR24993031
Ғылыми -- Техникалық Бағдарламасы бойынша «Табиғи антиоксиданттармен
және биологиялық белсенді заттармен байытылған күнделікті рационға
арналған пайдалы тағам өнімдерін дайындау технологиясын әзірлеу»
бағдарламалық-нысаналы қаржыландыру (ПЦФ) аясында қаржыландырылды.}

{\bfseries Әдебиеттер}

1. Krajewska A., Dziki D. Entrichment of cookies with fruits and their
by-products: chemical composition, andioxidant properties, and sensory
changes//Molecules. -2023. -28(10): 4005.
\href{https://doi.org/10.3390/molecules28104005}{DOI
10.3390/molecules28104005}.

2. Raszkowska E., et al. The use of blackcurrant pomace and erythritol
to optimase the functional properties of shortbread cookies//Scientific
Reports. -2024. -Vol.14: 3788.
\href{https://doi.org/10/1038/s41598-024-54461-7}{DOI
10/1038/s41598-024-54461-7}.

3. De Santis D., Ferri S., Rossi A., Frisoni R., Modesti M. Shortbread
cookies reformulation by raspberry powder enrichment: functional and
sensory aspects// International Journal of Food Science and Technology.
-2024. -Vol.59 (10). - P.7560-7569.
\href{https://doi.org/10/1111/ijfs.17526}{DOI 10/1111/ijfs.17526}.

5. Nour V., Blejan A.M, Codina G.G. Use of bilberry and blackcurrant
pomace powders as

Blackcurrant pomace powders as functional ingrediens in cookies //
Applied Science. -2025. -Vol.15(10): 5247. DOI 10.3390/app15105247.

6. Michalska A., et.al. Effect of different drying techniques on
physical properties of blackcurrant pomace powders // LWT Food Sci
Technol. -2017. -Vol.77.-P.114-121.
\href{https://doi.org/10/1016/j.lwt.2016.12.008}{DOI
10/1016/j.lwt.2016.12.008}.

7. Sadowskaa A, et.al. Properties and microstructure of blackcurrant
powders prepared using a new method of fluidized-bed jet milling and
drying versus other drying methods//CyTA - Journal of Food. -2019. -Vol.
17(1). -P.439-446.
\href{https://doi.org/10.1080/19476337.2019.1596985}{DOI10.1080/19476337.2019.1596985}.

8. Azman E.M., et.al. Effect of dehydrations on phenolic compounds and
antioxidant activity of blackcurrant (Ribes nigrum L.) pomace
//International Journal of Food Science and Technology. -2021.
-Vol.56(2). - P.600-607. \href{https://doi.org/10.1111/ijfs.14762}{DOI
10.1111/ijfs.14762}.

9. Mildner- Szkudlarz S., et.al. Physical and bio Ctive properties of
muffins enriched with raspberry and cranberry pomace powder: A Promising
Application of Fruit By-Products Rich in Biocompounds //Plant Foods for
Human Nutrition. -2016. -71. -P.165-173.
\href{https://doi/org/10/1007/s11130-016-0539-4}{DOI
10/1007/s11130-016-0539-4}.

10. Gornas P., et.al. The impact of different baking conditions on the
stability of extractable polyphenols in muffins enriched byfruit
pomace//LWT. -2016. -Vol.65 P.946-953.
\href{https://doi.org/10/1016/j/lwt.2015.09.029}{DOI
10/1016/j/lwt.2015.09.029}.

11. Lau K.Q., et al. Utilization of vegetable and fruit by-products
asfunctional ingredientsin foods: a review //Frontiers in Nutrition.
-2021. -Vol.8:661693. DOI
\href{https://doi/org/10.3389/fnut.2021.661693}{10.3389/fnut.2021.661693}.

12. Kumar H., et.al. Selected fruit pomaces: nutritional profile, health
benefits, and applications in functional foods and feeds //Current
Research in Food Science. -2024. -Vol.9: 100791.
\href{https://doi.org/?10.1016/j.crfs.2024.100791}{DOI
10.1016/j.crfs.2024.100791}.

13. Федина П.А., Яшин А.Я., Черноусова Н.И. Определение антиоксидантов в
продуктач растительного происхождения амперометрическим методом. //Химия
растительного сырья - 2010. - № 2. -С.91-97.

14. Idouraine A., Kohlhepp E.A., Weber C.W. Nutrient composition of
eight lines of naked seed pumkin (Cucurbita pepo L.)// Journal of
Agricultural and Food Chemistry. -1996. -Vol.44. P.721-724. DOI
\href{https://doi.org/10.1021/jf950630y}{10.1021/jf950630y}.

15. Kim M.Y., Kim E.J, Kim Y.N., Choi C., Lee B-H. Comparison of
chemical compositions and nutritive values of varions pumkin species and
parts // Nutrition research and practice. -2012. - 6(1). -P.21-27.
\href{https://doi.org/10.4162/nrp.2012.6.1.21}{DOI
10.4162/nrp.2012.6.1.21}.

16. Brzezowska J., Wojdylo A., Nowicka P., Turkiewicz I.P., Tracz K.,
Michalsla Ciechanowska A. How do the phenolics and amino acids affect
the antioxidant, antidiabetic and antiglycation properties of selected
fruit powders? //LWT. -2023. -Vol.189: 115454.
\href{https://doi.org/10/1016/j.lwt.2023.115454}{DOI
10.1016/j.lwt.2023.115454} .

{\bfseries References}

1. Krajewska A., Dziki D. Entrichment of cookies with fruits and their
by-products: chemical composition, andioxidant properties, and sensory
changes//Molecules. -2023. -28(10): 4005.
\href{https://doi.org/10.3390/molecules28104005}{DOI
10.3390/molecules28104005}.

2. Raszkowska E., et al. The use of blackcurrant pomace and erythritol
to optimase the functional properties of shortbread cookies//Scientific
Reports. -2024. -Vol.14: 3788.
\href{https://doi.org/10/1038/s41598-024-54461-7}{DOI
10/1038/s41598-024-54461-7}.

3. De Santis D., Ferri S., Rossi A., Frisoni R., Modesti M. Shortbread
cookies reformulation by raspberry powder enrichment: functional and
sensory aspects// International Journal of Food Science and Technology.
-2024. -Vol.59 (10). --P.7560-7569.
\href{https://doi.org/10/1111/ijfs.17526}{DOI 10/1111/ijfs.17526}.

5. Nour V., Blejan A.M, Codina G.G. Use of bilberry and blackcurrant
pomace powders as

Blackcurrant pomace powders as functional ingrediens in cookies //
Applied Science. -2025. -Vol.15(10): 5247. DOI 10.3390/app15105247.

6. Michalska A., et.al. Effect of different drying techniques on
physical properties of blackcurrant pomace powders // LWT Food Sci
Technol. -2017. -Vol.77. -P.114-121.
\href{https://doi.org/10/1016/j.lwt.2016.12.008}{DOI
10/1016/j.lwt.2016.12.008}.

7. Sadowskaa A, et.al. Properties and microstructure of blackcurrant
powders prepared using a new method of fluidized-bed jet milling and
drying versus other drying methods//CyTA - Journal of Food. -2019. -Vol.
17(1). -P.439-446.
\href{https://doi.org/10.1080/19476337.2019.1596985}{DOI
10.1080/19476337.2019.1596985}.

8. Azman E.M., et.al. Effect of dehydrations on phenolic compounds and
antioxidant activity of blackcurrant (Ribes nigrum L.) pomace
//International Journal of Food Science and Technology. -2021.
-Vol.56(2). - P.600-607. \href{https://doi.org/10.1111/ijfs.14762}{DOI
10.1111/ijfs.14762}.

9. Mildner- Szkudlarz S., et.al. Physical and bio Ctive properties of
muffins enriched with raspberry and cranberry pomace powder: A Promising
Application of Fruit By-Products Rich in Biocompounds //Plant Foods for
Human Nutrition. -2016. -71. -P.165-173.
\href{https://doi/org/10/1007/s11130-016-0539-4}{DOI
10/1007/s11130-016-0539-4}.

10. Gornas P., et.al. The impact of different baking conditions on the
stability of extractable polyphenols in muffins enriched byfruit
pomace//LWT. -2016. -Vol.65 P.946-953.
\href{https://doi.org/10/1016/j/lwt.2015.09.029}{DOI
10/1016/j/lwt.2015.09.029}.

11. Lau K.Q., et al. Utilization of vegetable and fruit by-products
asfunctional ingredientsin foods:a review //Frontiers in Nutrition.
-2021. -Vol.8:661693. DOI
\href{https://doi/org/10.3389/fnut.2021.661693}{10.3389/fnut.2021.661693}.

12. Kumar H., et.al. Selected fruit pomaces: nutritional profile, health
benefits, and applications in functional foods and feeds //Current
Research in Food Science. -2024. -Vol.9: 100791.
\href{https://doi.org/?10.1016/j.crfs.2024.100791}{DOI
10.1016/j.crfs.2024.100791}.

13. Fedina P.A., Jashin A.Ja., Chernousova N.I. Opredelenie
antioksidantov v produktach rastitel' nogo
proishozhdenija amperometricheskim metodom. //Himija
rastitel' nogo syr' ja - 2010. - № 2.
-S.91-97 {[}in Russian{]}

14. Idouraine A., Kohlhepp E.A., Weber C.W. Nutrient composition of
eight lines of naked seed pumkin (Cucurbita pepo L.)// Journal of
Agricultural and Food Chemistry. -1996. -Vol.44. P.721-724. DOI
\href{https://doi.org/10.1021/jf950630y}{10.1021/jf950630y}.

15. Kim M.Y., Kim E.J, Kim Y.N., Choi C., Lee B-H. Comparison of
chemical compositions and nutritive values of varions pumkin species and
parts // Nutrition research and practice. -2012. - 6(1). -P.21-27.
\href{https://doi.org/10.4162/nrp.2012.6.1.21}{DOI
10.4162/nrp.2012.6.1.21}.

16. Brzezowska J., Wojdylo A., Nowicka P., Turkiewicz I.P., Tracz K.,
Michalsla Ciechanowska A. How do the phenolics and amino acids affect
the antioxidant, antidiabetic and antiglycation properties of selected
fruit powders? //LWT. -2023. -Vol.189: 115454.
\href{https://doi.org/10/1016/j.lwt.2023.115454}{DOI
10.1016/j.lwt.2023.115454} .

\emph{{\bfseries Сведения об авторах}}

Омаралиева А.М. - кандидат технических наук, ассоциированный профессор,
Казахский университет технологии и бизнеса им. К. Кулажанова, Астана,
Казахстан, e-mail:
aigul-omar@mail.ru;

Ботбаева Ж.Т. - кандидат биологических наук, ассоциированный профессор,
Казахский университет технологии и бизнеса им. К. Кулажанова, Астана,
Казахстан, e-mail:
zhanar.b.t88@gmail.com;

Хастаева А.Ж. - PhD, ассоциированный профессор Казахский университет
технологии и бизнеса им. К. Кулажанова, Астана, Казахстан, e-mail:
gera\_or@mail.ru;

Бектурганова А.А. - кандидат технических наук, ассоциированный
профессор, Казахский университет технологии и бизнеса им. К. Кулажанова,
Астана, Казахстан, e-mail:
1968al1@mail.ru;

Агедилова М.Т. - кандидат химических наук, ассоциированный профессор,
Казахский университет технологии и бизнеса им. К. Кулажанова, Астана,
Казахстан, e-mail:
amt5@yandex.ru;

Рахимжанова А.Р. - магистр технических наук, Казахский университет
технологии и бизнеса им. К. Кулажанова, Астана, Казахстан, e-mail:
r.ayagoz@mail.ru;

Аубакирова А.К. - магистр технических наук, Казахский университет
технологии и бизнеса им. К. Кулажанова, Астана, Казахстан, e-mail:
alfiyakabylovna@gmail.com.

\emph{{\bfseries Information about the authors}}

Omaraliyeva A.M. - candidate of technical sciences, Associate Professor,
K. Kulazhanov Kazakh University of Technology and Business, Astana,
Kazakhstan, e-mail:
aigul-omar@mail.ru;

Botbayeva Zh.T. - candidate of technical sciences Associate Professor,
K. Kulazhanov Kazakh University of Technology and Business, Astana,
Kazakhstan, e-mail:
zhanar.b.t88@gmail.com;

Khastayeva A.Zh. - PhD, associate Professor, K. Kulazhanov Kazakh
University of Technology and Business, Astana, Kazakhstan, e-mail:
gera\_or@mail.ru;

Bekturganova A.A. - candidate of technical sciences Associate Professor,
K. Kulazhanov Kazakh University of Technology and Business, Astana,
Kazakhstan, e-mail:
1968al1@mail.ru;

Agedilova M.T. - candidate of chemical sciences, Associate Professor, K.
Kulazhanov Kazakh University of Technology and Business, Astana,
Kazakhstan, e-mail:
amt5@yandex.ru;

Rakhimzhanova A.R. - master of technical sciences, K. Kulazhanov Kazakh
University of Technology and Business, Astana, Kazakhstan, e-mail:
r.ayagoz@mail.ru;

Aubakirova A.K. - master of technical sciences, K. Kulazhanov Kazakh
University of Technology and Business, Astana, Kazakhstan, e-mail:
alfiyakabylovna@gmail.com.\
